\section{Related Work}\label{S6-3}

Our work sits at the intersection of three research streams: deep multiple instance learning for computational histopathology, risk-aware learning from imbalanced medical data, and trustworthy clinical AI via selective classification. We review each stream below, highlighting why existing approaches, when applied in isolation, fall short of delivering clinically safe MIP detection---and how their integration gives rise to the framework proposed in this paper.

\subsection{Deep MIL for Histopathology Subtyping}

Deep learning has transformed computational pathology by enabling automated analysis of gigapixel whole-slide images (WSIs) \cite{coudray2018classification, campanella2019clinical}. Because WSIs far exceed the input dimensions of standard neural networks, multiple instance learning (MIL) has become the dominant paradigm: each slide is treated as a bag of image tiles, and tile-level features are aggregated into a slide-level prediction without requiring pixel-level annotations \cite{ilse2018attention}. Early MIL methods relied on attention mechanisms to highlight diagnostically relevant regions, and subsequent frameworks such as CLAM \cite{lu2021data} introduced clustering-based regularization to improve sample efficiency. More recently, the backbone has shifted from convolutional networks to Vision Transformers (ViTs) \cite{dosovitskiy2020image} combined with self-supervised pre-training. Histology-specific foundation models such as Phikon \cite{filiot2023scaling}, pre-trained via masked image modeling on large-scale pathology corpora, now deliver substantially stronger feature representations than generic ImageNet initialization. Transformer-based MIL architectures have further been tailored specifically for lung cancer image classification \cite{an2024transformer}.

Despite these advances, the evaluation of MIL models for pathology subtyping---including prior work on micropapillary (MIP) detection---has focused predominantly on aggregate metrics such as AUC and overall accuracy. The clinically asymmetric cost of errors is largely ignored. In high-risk settings such as MIP lung adenocarcinoma, where even a minor component ($>5\%$) carries severe prognostic implications \cite{caso2020underlying, dai2017tumor}, a model optimized solely for average performance may still produce an unacceptably high false-negative rate. Our work begins from a strong ViT-MIL baseline but explicitly identifies and addresses this neglected safety dimension.

\subsection{Risk-Aware Learning: From Cost-Sensitive Reweighting to Neyman--Pearson Guarantees}

Standard classification objectives assume symmetric error costs and balanced classes, which poorly align with medical diagnosis where rare, high-risk conditions demand stricter control of false negatives \cite{he2009learning}. Cost-sensitive learning (CSL) partially mitigates this by re-weighting the loss to penalize false negatives more heavily than false positives \cite{elkan2001foundations}. A recent comprehensive review of 173 studies confirms that CSL is among the most widely adopted strategies for handling imbalanced medical data, with applications spanning radiology, pathology, and clinical risk prediction \cite{araf2024cost}. However, CSL suffers from a fundamental limitation: the class weights are heuristically tuned, and the resulting model provides no formal, distribution-free guarantee on the false-negative rate (FNR). The achievable FNR depends on the chosen weight, the data distribution, and the model architecture in ways that are difficult to characterize a priori.

A more rigorous alternative is the Neyman--Pearson (NP) classification framework, which formulates learning as a constrained optimization problem: minimize the type-II error (FNR) while bounding the type-I error (FPR) below a user-specified level \cite{rigollet2011neyman}. The NP umbrella algorithm \cite{tong2018neyman} constructs a threshold via order statistics on a held-out positive calibration set, yielding a finite-sample, distribution-free guarantee that the FNR remains below the target level with high probability. Recent work has extended NP classification to multi-class settings through cost-sensitive reformulations \cite{tian2025neyman}, underscoring its relevance for safety-critical applications.

Yet NP calibration, when applied to a naively trained model, comes at a steep practical cost: enforcing a stringent FNR target (e.g., $\leq 1\%$) forces the decision threshold so low that specificity collapses. Moreover, NP methods require holding out positive samples for calibration, which further reduces the already scarce minority-class training data. Critically, existing work treats CSL and NP as independent strategies, without exploring whether CSL can serve as an effective training precursor to NP calibration. Our framework explicitly investigates this synergy: CSL shifts the positive-class score distribution upward during training, enabling NP calibration to meet the same FNR target at a less aggressive threshold, thereby preserving specificity.

\subsection{Trustworthy Clinical AI: Selective Classification and Conformal Risk Control}

Even with rigorous FNR control at the decision-threshold level, a deployed model will encounter ambiguous or out-of-distribution samples on which its predictions should not be trusted. Selective classification addresses this by endowing models with a reject option: the classifier may abstain from making a prediction, deferring uncertain cases to human experts \cite{el2010foundations}. Architecturally, SelectiveNet \cite{geifman2019selectivenet, geifman2017selective} jointly trains a prediction head and a selection head, optimizing predictive performance on the covered (accepted) subset while satisfying a target coverage constraint via a differentiable penalty. This approach has been extended to handle distribution shifts and hierarchical rejection schemes \cite{goren2024hierarchical}. In clinical imaging, the value of selective deferral has been demonstrated in breast cancer screening \cite{leibig2022combining} and in general AI-assisted diagnosis pipelines \cite{dvijotham2023enhancing}.

Determining the appropriate rejection threshold is a post-training calibration problem. The conformal risk control (CRC) framework \cite{angelopoulos2022conformal, bates2021distribution} provides a general paradigm for this task: given a pre-trained model, a held-out calibration set, and a monotone risk function of interest, one selects the most permissive threshold such that the empirical risk remains below a user-specified target level $\alpha$. This calibration strategy requires no retraining, accommodates arbitrary risk metrics, and serves as a principled alternative to fixed heuristic thresholds. In this work, we adopt a CRC-style calibration procedure: after training a SelectiveNet backbone, we determine the rejection threshold $\tau^*$ on a held-out calibration set by enforcing $\widehat{\mathrm{FNR}}(\tau) \leq \alpha$ or $\hat{R}(\tau) \leq \alpha$, depending on the clinical objective. The formal finite-sample guarantees on the resulting FNR are provided by our Neyman--Pearson umbrella procedure (Theorem~1) and the risk--coverage tradeoff (Theorem~2), which complement the CRC calibration paradigm.

Critically, existing selective classification methods---whether based on SelectiveNet or CRC calibration---control \emph{overall} selective risk (the average error rate on accepted samples). In asymmetric-error medical scenarios such as MIP detection, this is insufficient: a model may satisfy an overall risk target of $5\%$ while still missing over $12\%$ of true MIP cases, as we show in our experiments. The class-conditional false-negative rate is not directly constrained by standard selective risk minimization. To our knowledge, no prior work has integrated cost-sensitive training objectives with FNR-aware selective classification to explicitly bound missed diagnoses on the accepted subset. Our framework addresses this gap by introducing four complementary operating cases that span the full spectrum from general risk control to strict, class-conditional FNR control, with and without cost-sensitive training, allowing clinicians to select the safety operating point most appropriate for their deployment context.
